Let $\mathcal D$ be the set of Dirichlet pairs (\cref{I.adele.def:dirichlet-pair}) and $w\mu=(\mu_2,\mu_1)$ for $\mu=(\mu_1,\mu_2)\in\mathcal D$. An \emph{orbit} is a subset
$o\subseteq\mathcal D$ of the form $o=\{\mu,w\mu\}$ with $\mu\in\mathcal D$ (as in \cref{I.spec.thm:cuspidal-data-orthogonality}); it has one element if
$\mu_1=\mu_2$ and two otherwise. Let $\mathfrak O$ be the set of orbits. The set of \emph{cuspidal data} is the disjoint union
\[
  \mathfrak X=\{\mathrm{cusp}\}\sqcup\mathfrak O ,
\]
where $\mathrm{cusp}$ is a single symbol, the datum that stands for all cuspidal data $(G,\sigma)$ together.
